\documentclass[abstracton]{scrartcl}

\usepackage[a4paper, total={15.5cm, 23cm}]{geometry}


\usepackage{amsmath}
\usepackage{amsthm}
\usepackage{amssymb}

\newtheorem{theorem}{Theorem}
\newtheorem{lemma}[theorem]{Lemma}
\newtheorem{cor}[theorem]{Corollary}
\newtheorem{proposition}[theorem]{Proposition}
\newtheorem{definition}[theorem]{Definition}
\newtheorem{example}[theorem]{Example}

\newcommand{\cU}{\mathcal{U}}
\newcommand{\X}{\mathcal{X}}
\newcommand{\conv}{\textup{conv}}

\title{Remember LP}
\date{}

\begin{document}
	
\maketitle
	
Ursprüngliches Problem:

\begin{align}
	\min\; & t \label{eq:mirror-obj}\\[1mm]
	\text{s.t.}\quad
	& \max_{s\in[k]} \sum_{i\in\mathcal{S}(w)} c_i^s
	\;\le\;
	t \sum_{i\in\mathcal{S}(w)} w_i,
	\label{eq:mirror-i}
	\\[1mm]
	& \sum_{i\in\mathcal{S}(w)} w_i
	\;\le\;
	\mathrm{OPT}_{\mathrm{dual}} .
	\label{eq:mirror-ii}
\end{align}
	
Bedingungen ersetzt durch sufficient constraints:
\begin{alignat}{3}
	\max \quad & t \\[1mm]
	\text{s.t.}\quad
	& t \sum_{i\in\mathcal S} c_i^s \;\le\; \sum_{i\in\mathcal S} w_i
	&& \forall s\in[k],\ \forall \mathcal S\subseteq[n]:|\mathcal S|=p, \label{eq:pd-lp-i}\\[1mm]
	& \sum_{i\in\mathcal{S}(w)} w_i \;\le\; p\,\alpha - \sum_{i=1}^n \gamma_i, \label{eq:pd-lp-ii}\\[1mm]
	& \alpha - \sum_{s=1}^k c_i^s \beta_s \;\le\; \gamma_i
	&& \forall i\in[n], \label{eq:pd-lp-dualfeas}\\[1mm]
	& w_i \;=\; \sum_{s=1}^k \beta_s c_i^s
	&& \forall i\in[n], \label{eq:pd-lp-w}\\[1mm]
	& \sum_{s=1}^k \beta_s = 1,\quad \beta_s \ge 0
	&& \forall s\in[k], \label{eq:pd-lp-simplex}\\[2mm]
	& \gamma_i \ge 0 && \forall i\in[n], \qquad \alpha\in\mathbb{R}. \label{eq:pd-lp-sign}
\end{alignat}

Nach Dualem Ersatz und zeigen, dass 6 \& 7 immer erfüllt werden können:

\begin{alignat}{3}
	\label{lp:compact}
	\max \quad & t \\[1mm]
	\text{s.t.}\quad
	& w_i \;=\; \sum_{s=1}^k \beta_s\, c_i^s
	&& \forall i\in[n], \label{lp:compact-w}\\[1mm]
	& \sum_{s=1}^k \beta_s \;=\; 1 \label{lp:compact-simplex}\\[2mm]
	& p\,\alpha_s \;-\; \sum_{i=1}^n \lambda^s_i \;\ge\; 0
	&& \forall s\in[k], \label{lp:compact-dual1}\\[1mm]
	& \lambda^s_i \;\ge\; \alpha_s - w_i + t\,c_i^s
	&& \forall s\in[k],\ \forall i\in[n], \label{lp:compact-dual2}\\[1mm]
	& \lambda^s_i \ge 0
	&& \forall s\in[k],\ \forall i\in[n], \label{lp:compact-dual3}\\[1mm]
	& t \ge 0,\quad \alpha_s \in \mathbb{R},\quad \beta_s \ge 0
	&& \forall s\in[k]. \label{lp:compact-sign}
\end{alignat}

"Remember" LP (Teil der Items fixiert):

\begin{alignat}{3}
	\label{lp:compact-remember}
	\max \quad & t \\[1mm]
	\text{s.t.}\quad
	& w_i \;=\; \sum_{s=1}^k \beta_s\, c_i^s
	&& \forall i\in R, \label{lp:remember-w}\\[1mm]
	& \sum_{s=1}^k \beta_s \;=\; 1
	&& \label{lp:remember-simplex}\\[2mm]
	& p_{\mathrm{rem}}\,\alpha_s \;-\; \sum_{i\in R} \lambda^s_i
	\;+\; \sum_{s=1}^k \beta_s \sum_{i\in F} c_i^s
	\;-\; t \sum_{i\in F} c_i^s \;\ge\; 0
	&& \forall s\in[k], \label{lp:remember-dual1}\\[1mm]
	& \lambda^s_i \;\ge\; \alpha_s - w_i + t\,c_i^s
	&& \forall s\in[k],\ \forall i\in R, \label{lp:remember-dual2}\\[1mm]
	& \lambda^s_i \;\ge\; 0
	&& \forall s\in[k],\ \forall i\in R, \label{lp:remember-dual3}\\[1mm]
	& t \;\ge\; 0,\quad \alpha_s \in \mathbb{R},\quad \beta_s \;\ge\; 0
	&& \forall s\in[k]. \label{lp:remember-sign}
\end{alignat}

Herleitung:

\begin{itemize}
	\item Menge $F$ von Items fix gewählt; wir müssen nur noch
	$p_{\mathrm{rem}} = p - |F|$ weitere Items aus  Restpool $R$ wählen
	
	\item Die Bedingung (5) soll jetzt für die vollständige Menge 
	$F \cup T$ ($T \subseteq R$)gelten
	\[
	t \sum_{i\in F\cup T} c_i^s \;\le\; \sum_{i\in F\cup T} w_i.
	\]
	
	\item Festen und variablen Teil trennen
	\[
	\sum_{i\in T} (w_i - t c_i^s)
	\;\ge\;
	- \sum_{i\in F} (w_i - t c_i^s).
	\]
	
	\item Dann definieren wir
	\[
	d_i^s := w_i - t c_i^s \qquad (i\in R).
	\]
	Dann lautet die Bedingung äquivalent:
	\[
	\sum_{i\in T} d_i^s \;\ge\; 
	 - \sum_{i\in F}(w_i - t c_i^s)
	\qquad\forall T\subseteq R,\ |T| = p_{\mathrm{rem}}.
	\]
	
	\item wir betrachten dieses selection-problem:
	\[
	\min\Bigl\{\sum_{i\in R} d_i^s x_i : 
	\sum_{i\in R} x_i = p_{\mathrm{rem}},\ x_i\ge 0 \Bigr\}.
	\]
	TU, daher entspricht der Optimalwert exakt
	dem diskreten Minimum.
	
	\item Das Dual:
	\[
	\max\Bigl\{
	p_{\mathrm{rem}}\alpha_s - \sum_{i\in R} \lambda_i^s :
	\alpha_s - \lambda_i^s \ge d_i^s,\ \lambda_i^s\ge 0
	\Bigr\}.
	\]
	
	\item Nach starker Dualität gilt die Bedingung genau dann, wenn
	der Dualwert mindestens $- \sum_{i\in F}(w_i - t c_i^s)$ ist:
	\[
	p_{\mathrm{rem}}\alpha_s - \sum_{i\in R} \lambda_i^s \;\ge\; - \sum_{i\in F}(w_i - t c_i^s).
	\]
	
	\item also
	\[
	p_{\mathrm{rem}}\alpha_s - \sum_{i\in R} \lambda_i^s
	+ \sum_{i\in F} w_i - t \sum_{i\in F} c_i^s \;\ge\; 0.
	\]
	
	\item Da für alle Items (insbesondere für $i\in F$) gilt
	\[
	w_i = \sum_{s=1}^k \beta_s c_i^s,
	\]
	können wir die Summe über feste Items ausdrücken als
	\[
	\sum_{i\in F} w_i
	= \sum_{s=1}^k \beta_s \sum_{i\in F} c_i^s.
	\]
	
	\item Damit ergibt sich die Remember-Bedingung
	\[
	p_{\mathrm{rem}} \alpha_s
	- \sum_{i\in R} \lambda_i^s
	+ \sum_{s=1}^k \beta_s \sum_{i\in F} c_i^s
	- t \sum_{i\in F} c_i^s
	\;\ge\; 0.
	\]
	
	\item Da $\sum_{i\in F} w_i = \sum_{s=1}^k \beta_s \sum_{i\in F} c_i^s$
	gilt, müssen wir für feste Items keine eigenen Variablen $w_i$ einführen.
	
	\item Alle übrigen Bedingungen des LP bleiben unverändert;. Nur die Indexmengen werden von $[n]$ auf $R$ reduziert.
\end{itemize}





\end{document}